\begin{center}
{\fontsize{14}{17}\selectfont\bfseries Introducción}
\end{center}
\addcontentsline{toc}{section}{Introducción}

La rápida adopción de la Inteligencia Artificial (IA) en entornos organizacionales ha generado una demanda creciente de sistemas inteligentes capaces de operar eficazmente sobre bases de conocimiento internas y específicas del dominio. En contextos industriales donde los datos operativos confidenciales, los procesos propietarios, la documentación regulatoria y el conocimiento institucional no están disponibles públicamente. Los Modelos de Lenguaje Grandes (LLM, \textit{Large Language Models}) de propósito general enfrentan una limitación estructural crítica: son entrenados sobre entornos de la nube y en algunos casos se ha visto el uso de manera pública información privada, por tanto existen brechas de seguridad o se torna complicado el acceso inherente a la información privada para una fácil retroalimentación, misma que se califica como sensible para la organización y que es la que impulsa la toma de decisiones cotidiana.

RAG \textit{(Retrieval Augmented Generation)} mejor conocido en español como  Generación Aumentada por Recuperación conocida, ha ganado atención como mecanismo para anclar las respuestas de los modelos de lenguaje en documentos relevantes y actualizados recuperados de repositorios internos. Las arquitecturas RAG disocian el conocimiento paramétrico incrustado en los pesos del modelo del conocimiento fáctico almacenado en fuentes externas, permitiendo que los modelos generen respuestas informadas por contenido específico de la organización sin requerir reentrenamiento o ajuste fino costosos (\cite{Lewis2020rag}). Aun cuando RAG extiende considerablemente el alcance funcional de los modelos de IA hacia dominios de conocimiento privados, introduce una nueva clase de problemas: el modelo puede recuperar información irrelevante o incompleta, no reconocer los límites de su propio conocimiento, o producir respuestas confiadas pero incorrectas; fenómeno conocido comúnmente como alucinación. Complementario al RAG, el mecanismo de autorreflexión (la capacidad de un modelo de lenguaje para introspeccionar sobre su propio razonamiento, identificar inconsistencias y refinar iterativamente sus respuestas), ha emergido como una dirección prometedora para mejorar la confiabilidad de las respuestas y la calibración en la capacidad de respuesta. Los agentes autorreflexivos, como los descritos en el marco Reflexion (\cite{Shinn2023_Reflexion}) o mediante las arquitecturas Self-RAG (\cite{Asai2024_Selfrag}), son capaces de evaluar la calidad de la evidencia recuperada, detectar brechas lógicas en las respuestas generadas y desencadenar ciclos correctivos de recuperación o razonamiento. Estas capacidades son particularmente valiosas en entornos empresariales donde las respuestas incorrectas o mal fundamentadas pueden tener consecuencias operativas, legales o financieras.

Un desafío poco explorado concierne al despliegue de estos mecanismos en modelos de lenguaje pequeños con recursos limitados y orientados a dominios específicos. Si bien los LLM de frontera (ej. GPT-4, Claude 3) poseen capacidad paramétrica suficiente para implementar comportamientos complejos de razonamiento y reflexión, las organizaciones frecuentemente operan bajo restricciones de recursos (\textit{hardware}) que impiden el despliegue de modelos muy grandes. Los Modelos de Lenguaje Pequeños (SLM, \textit{Small Language Models}), en el rango de 1B a 13B (1.000 a 13.000 millones de parámetros), son más viables para el despliegue en infraestructuras propias, con preservación de la privacidad y eficiencia de costos. No obstante, si estos modelos más pequeños pueden inducir eficazmente comportamientos autorreflexivos y aprovechar pipelines RAG con dominio específico, en el caso particular el soporte de conocimiento interno empresaial; sigue siendo una pregunta abierta y crítica y constituye el núcleo del problema que se aborda mas adelante.

Aún con el creciente número de trabajos sobre arquitecturas RAG y agentes LLM autorreflexivos con una tendencia en la implementación de soluciones de este tipo, la intersección de estos mecanismos en el contexto de modelos de lenguaje pequeños o entornos con recursos limitados, orientados al conocimiento empresarial interno, sigue siendo un campo insuficientemente explorado. Los estudios existentes centran principalmente hallazgos sobre modelos de gran escala en contextos generales o en RAG como problema de recuperación de información, sin considerar las restricciones específicas de los despliegues industriales. Se realizó un estudio de mapeo sistemático SMS \textit{(Systematic Mapping Study)}(\cite{Alcides2026}), con el propósito de analizar la literatura existente sobre trabajos que emplean o integren componente de RAG y mecanismos de autorreflexión en modelos de lenguaje aplicados al soporte de conocimiento. revisando y explorando diversos enfoques propuestos, con énfasis en aquellas aplicables a tareas de conocimiento específicas del dominio o de entornos empresariales  sobre la cual se fundamenta la presente propuesta doctoral. A partir de esta base, el presente trabajo se centra en el diseño, la implementación y la evaluación de una arquitectura de autorreflexión integrada a un pipeline RAG sobre modelos de lenguaje pequeños (SLM) aplicada al soporte de conocimiento empresarial interno en procesos de inducción de personal.

\newpage
